\chapter{Recetas paso a paso}
\label{cap:recetas}

Cada receta es un cambio típico, con los archivos exactos y la forma de comprobarlo. Si una
receta de aquí y un \emph{runbook} (\archivo{docs/runbooks/}) no coinciden, manda el runbook:
es el procedimiento vigente; esta guía lo explica con ejemplos.

\begin{idea}[Antes de cualquier receta]
\begin{enumerate}
  \item Crea una rama: \cod{git switch -c feature/nombre} (nunca se trabaja en \cod{master}).
  \item Abre la ficha del área que vas a tocar (\archivo{docs/areas/}).
  \item Si hay una trampa para esa zona, léela (capítulo~\ref{cap:trampas}).
  \item Al terminar: \cod{cargo build --release}, \cod{cargo test} y
  \cod{python3 scripts/docs.py check}.
\end{enumerate}
\end{idea}

%% ------------------------------------------------------------------------------------------
\section{Añadir una opción}

\emph{Ejemplo:} «Mostrar segundos en la barra», con valores sí o no. Es la receta más común y
\emph{no toca Rust}.

\begin{enumerate}
  \item \textbf{La fila.} En \archivo{qml/SettingsSheet.qml}, añade un objeto a
  \cod{optionRows}. \cod{cat} es la categoría de la barra lateral (\cod{reproduccion},
  \cod{apariencia} o \cod{vistas}); \cod{fallback} es el valor por defecto.
\begin{qml}
{ cat: "vistas", key: "segundos", label: "Mostrar segundos", fallback: "si",
  options: [{ id: "si", label: "Sí" }, { id: "no", label: "No" }] },
\end{qml}
  Con eso la fila aparece y se guarda sola en \cod{settings.json}.
  \item \textbf{Leerla.} Donde la uses, con \cod{void AppSettings.revision} delante: sin eso,
  el enlace no se vuelve a evaluar al cambiar la opción.
\begin{qml}
readonly property bool segundos: { void AppSettings.revision; return AppSettings.opcion("segundos") !== "no" }
\end{qml}
  Si la leen varios archivos, ponla en \archivo{Theme.qml}.
  \item \textbf{Documentarla.} Añade la clave al \cod{opciones:} de la ficha de su área y
  explica su efecto. \cod{docs.py check} \emph{falla} si una clave del código no está en
  ninguna ficha.
  \item \textbf{Comprobar.} Una captura de la categoría
  (\cod{--shot ruta.png --settings-cat vistas}) y otra con la opción puesta (un
  \cod{settings.json} de prueba en el \cod{HOME} temporal).
\end{enumerate}

\begin{trampa}[Errores habituales]
\begin{itemize}
  \item Olvidar \cod{void AppSettings.revision}: el valor no cambia hasta reiniciar.
  \item Llamar \cod{setOpcion(clave, "")}: borra la opción y deja el \cod{fallback}.
  \item Un botón de opción sin 44 px de área o sin foco visible.
\end{itemize}
\end{trampa}

%% ------------------------------------------------------------------------------------------
\section{Añadir un color a las skins}

\begin{enumerate}
  \item En \archivo{qml/Theme.qml}:
\begin{qml}
readonly property color miColor: pick("miColor", "#FF8800")
\end{qml}
  El segundo valor es el de por defecto (el de la skin Vidrio).
  \item Añade \cod{"miColor": "#..."} a la sección \cod{colores} de \emph{las cuatro} skins
  de \archivo{skins/} (\cod{vidrio}, \cod{medianoche}, \cod{grafito}, \cod{neon}).
  \item Úsalo como \cod{Theme.miColor}. Nunca un color escrito a mano en un componente.
  \item Si las personas que usan la app pueden ponerlo en su skin, añádelo a la lista de
  claves del README.
  \item Captura con al menos dos skins (Vidrio y Neón).
\end{enumerate}

%% ------------------------------------------------------------------------------------------
\section{Añadir un archivo QML}

\begin{enumerate}
  \item Crea \archivo{qml/MiCosa.qml}.
  \item Regístralo en \archivo{build.rs}, en \cod{QML_FILES} (o en \cod{QML_SINGLETONS} si
  es un singleton, con \cod{pragma Singleton}).
  \item Añádelo al \cod{archivos:} de la ficha de su área y a su tabla «Archivos».
  \item Úsalo como \cod{MiCosa { ... }} desde cualquier QML.
\end{enumerate}

Si olvidas el paso 2, QML dirá que el tipo no existe; \cod{docs.py check} también lo detecta.

%% ------------------------------------------------------------------------------------------
\section{Añadir una propiedad o una función a un singleton}

\emph{Ejemplo:} saber cuántas canciones quedan en la cola, con la propiedad \cod{restantes}.

\begin{enumerate}
  \item \textbf{La propiedad}, en el bloque \cod{extern "RustQt"} de \archivo{src/player.rs},
  junto a las demás:
\begin{rust}
#[qproperty(i32, restantes)]
\end{rust}
  \item \textbf{El campo}, en \cod{struct PlayerRust}, y su valor inicial en
  \cod{impl Default for PlayerRust}:
\begin{rust}
restantes: i32,          // en el struct
restantes: 0,            // en Default
\end{rust}
  \item \textbf{Actualizarla} donde cambie la cola o la pista; por ejemplo al final de
  \cod{load}:
\begin{rust}
let restantes = (self.tracks.len() - index - 1) as i32;
self.as_mut().set_restantes(restantes);
\end{rust}
  \item \textbf{Usarla en QML}, con el nombre en \cod{camelCase} (aquí ya coincide):
\begin{qml}
Text { text: Player.restantes + " canciones más" }
\end{qml}
\end{enumerate}

Para una \emph{función}, declárala con \cod{#[qinvokable]} en el bloque del puente
(\emph{solo la firma}) y escribe el cuerpo en \cod{impl qobject::Player}. Si es lenta, ejecútala
en un hilo con \cod{qt_thread().queue(...)} y protégela con un contador (capítulo~\ref{cap:puente}).

\begin{trampa}
El cuerpo \emph{nunca} va dentro del bloque \cod{extern "RustQt"}: cxx-qt falla con
«Unsupported verbatim input in ForeignItem».
\end{trampa}

\subsection{Un singleton de Rust nuevo}

\begin{enumerate}
  \item \archivo{src/mi\_cosa.rs} con su \cod{#[cxx_qt::bridge]}.
  \item En \archivo{build.rs}, dentro de \cod{.files([...])}.
  \item En \archivo{src/main.rs}, \cod{mod mi_cosa;}.
  \item En el \cod{archivos:} de su ficha.
  \item \cod{cargo build --release} sin avisos.
\end{enumerate}

%% ------------------------------------------------------------------------------------------
\section{Añadir datos nuevos de YouTube Music}

\emph{Ejemplo:} una sección con las «Subidas» del usuario.

\begin{enumerate}
  \item \textbf{Descubre qué pide la web.} Abre music.youtube.com con las herramientas de
  desarrollador (pestaña Red), filtra por \cod{browse} y anota el \cod{browseId} de la petición.
  \item \textbf{Pruébalo sin interfaz.} \archivo{ytm.rs} ya tiene \cod{browse(browse_id)},
  que devuelve \cod{Vec<Zone>}. Añade un caso a \cod{probe} en \archivo{main.rs} y ejecuta
  \cod{--probe subidas}.
  \item \textbf{Publica en \cod{Library}.} Una propiedad (\cod{#[qproperty(QString, subidas)]})
  con su campo, y una entrada en el array \cod{sections} de \cod{load()}, como la de
  \cod{megusta}. El array tiene tamaño fijo: sube también el número de su tipo
  (\cod{[...; 5]} a \cod{[...; 6]}); el compilador te lo recordará.
  \item \textbf{Conviértelo en QML.} En \archivo{Catalog.qml}: una propiedad
  \cod{subidas: parse(Library.subidas)}, el \cod{case "subidas"} en \cod{zonas()} y la entrada
  en \cod{sections} (con \cod{cuenta: true} si exige sesión). La pestaña aparece sola en la
  navegación y en «Vista de cada sección».
  \item \textbf{La vista por defecto}, en \cod{DEFAULT_VIEWS} de \archivo{settings.rs}.
\end{enumerate}

%% ------------------------------------------------------------------------------------------
\section{Añadir una entrada al menú de clic derecho}

\emph{Ejemplo:} «Mover al final» en las canciones de la cola.

\begin{enumerate}
  \item En \archivo{ContextMenu.qml}, la propiedad \cod{entries} construye la lista según lo
  pulsado. Las pistas de la cola llegan con \cod{queueIndex} (lo añade \archivo{QueueList}),
  así que dentro de \cod{if (inQueue)}:
\begin{qml}
out.push({ id: "final", label: "Mover al final", icon: Theme.iconAddQueue })
\end{qml}
  \item En \cod{run(action)}, qué hace:
\begin{qml}
} else if (action === "final") {
    Player.moveInQueue(target.queueIndex, Catalog.queue.length - 1)
    notify("«" + target.title + "» al final de la cola")
\end{qml}
  \item Si la lista va a ser larga (como las playlists donde guardar), no la metas en el menú
  principal: haz una \emph{subpágina}. Una entrada con \cod{id: "page:loquesea"} y
  \cod{more: true} abre la página \cod{loquesea}, y \cod{entries} devuelve su contenido cuando
  \cod{page === "loquesea"}.
  \item Pruébalo con un clic simulado:
\begin{terminal}
MUSICPLAYER_DEBUG_CLICK="1224,300,right,6000" ... --shot /tmp/menu.png --autoplay --tab cola
\end{terminal}
\end{enumerate}

%% ------------------------------------------------------------------------------------------
\section{Añadir un atajo de teclado}

En \archivo{Main.qml}, junto a los demás:

\begin{qml}
Shortcut { sequence: "L"; enabled: !win.settingsOpen && !searchField.editing
           onActivated: win.lyricsOn = !win.lyricsOn }
\end{qml}

La condición \cod{!searchField.editing} evita que la tecla se «coma» cuando estás escribiendo
una búsqueda; \cod{!win.settingsOpen} hace lo mismo con Opciones.

%% ------------------------------------------------------------------------------------------
\section{Añadir un argumento de depuración}

Un argumento nuevo (para una captura) se añade a \archivo{Main.qml} con \cod{arg("--nombre")}
y a la tabla de argumentos de \archivo{docs/runbooks/verificar-interfaz.md}, \emph{en la misma
tarea}. Si es una sonda sin ventana, un caso nuevo en \cod{probe} de \archivo{main.rs} y su
fila en la tabla de \archivo{verificar-reproduccion.md}.

%% ------------------------------------------------------------------------------------------
\section{Crear una mascota}

No hace falta código: una carpeta en \cod{<config>/musicplayer/mascotas/<id>/} con un
\cod{mascota.json} y los PNG de las poses (formato en el capítulo~\ref{cap:datos}). Solo
\cod{reposo} es obligatoria. La lógica que elige la pose vive en \archivo{Mascot.qml}; si
quieres un estado nuevo, se añade ahí y en su tabla de sustitución.

\section{Crear una skin}

Tampoco necesita código: un \cod{<id>.json} (o una carpeta con \cod{skin.json} y una imagen de
fondo) en \cod{<config>/musicplayer/skins/}, con \cod{nombre} y \cod{colores}. Pulsa
«recargar» en Opciones. Para una skin que se incluya en el programa, copia el JSON a
\archivo{skins/} y añádelo a \cod{BUILTIN_SKINS} de \archivo{settings.rs}; ojo: un color nuevo
va en las cuatro.

%% ------------------------------------------------------------------------------------------
\section{Escribir una prueba}

Las pruebas están al final del archivo Rust, en un módulo \cod{#[cfg(test)]}:

\begin{rust}
#[cfg(test)]
mod tests {
    use super::*;

    #[test]
    fn quita_el_numero_de_reproducciones() {
        assert_eq!(artist_from_subtitle("Quevedo • 75 M reproducciones"), "Quevedo");
    }
}
\end{rust}

\begin{terminal}
cargo test --release
\end{terminal}

Es lo que da más seguridad por menos esfuerzo en lo que es \emph{parseo de datos ajenos}:
formatos de letra, firmas, convertir una respuesta de YouTube. Lo que habla con la red no tiene
pruebas (está en la deuda conocida); se prueba con las sondas.

\section{Añadir un plugin}

Un plugin nuevo (por ejemplo, otro servicio de «escuchando ahora») sigue el patrón de
\archivo{plugins/discord.rs}:

\begin{enumerate}
  \item Un archivo \archivo{src/plugins/mi\_plugin.rs} con un tipo que sabe conectarse y recibir
  la pista, la posición y la parada. Declara \cod{mod mi_plugin;} en \archivo{plugins.rs}.
  \item Su configuración, como campos nuevos de \cod{Config} (con \cod{#[serde(default)]}) para
  que \cod{plugins.json} de versiones anteriores siga siendo válido.
  \item En \cod{worker}, una variable \cod{Option<MiPlugin>} que se crea y se destruye según la
  configuración, y una rama por evento que le avisa.
  \item Las funciones invocables de \cod{Plugins} (\cod{set_mi_plugin}) y su propiedad de estado.
  \item La pestaña «Plugins» de \archivo{SettingsSheet.qml} y su ficha
  (\archivo{docs/areas/plugins.md}).
\end{enumerate}

Nunca llames a un plugin desde \cod{Player}: el hilo \cod{worker} existe para que la red de un
plugin no frene la interfaz (ADR-15).

\section{Añadir un área, una decisión o una trampa}

\begin{itemize}
  \item \textbf{Un área nueva:} \archivo{docs/runbooks/nueva-area.md} (copia
  \archivo{docs/templates/area.md}, rellena el frontmatter y ejecuta
  \cod{python3 scripts/docs.py index}).
  \item \textbf{Una decisión} que cruza áreas, o que fija el propietario:
  \archivo{docs/runbooks/registrar-decision.md}.
  \item \textbf{Una trampa} que costó encontrar: una línea en \archivo{docs/TRAMPAS.md}, con
  síntoma, causa y arreglo.
\end{itemize}

\section{Terminar una tarea}

Es el \emph{checklist final} de \archivo{docs/AGENTES.md} (§12), resumido:

\begin{enumerate}
  \item Nada lento en el hilo de Qt; las respuestas obsoletas se descartan.
  \item Colores por \cod{Theme}; opciones por \cod{AppSettings.opcion}, con su fila y en su ficha.
  \item QML y puentes nuevos registrados.
  \item \cod{cargo build --release} sin avisos, \cod{cargo test} en verde,
  \cod{python3 scripts/docs.py check} sin problemas.
  \item Captura de cada pantalla tocada; \cod{--autoplay} si tocaste reproducción; medida en
  release si cambia memoria o tiempos.
  \item Ficha del área al día, con su entrada en el historial; ADR o trampa si las hubo.
  \item Datos del propietario como estaban; sin procesos de prueba vivos.
  \item Dicho qué \emph{no} se verificó (Linux, Wayland, el sonido real, con sesión\dots).
\end{enumerate}

Y, solo con el visto bueno de quien manda en el proyecto, el commit y la unión a \cod{master}
(\cod{git merge --no-ff}, y borrar la rama).
